\par\Needspace{6\baselineskip}\medskip
\hypertarget{native-126-article093}{}\pdfbookmark[2]{93.}{article-093}
\begin{center}93.\end{center}
Ut solutionem aequationis 15 pro casibus in praxi frequentissimis quantum fieri potest commodissimam reddamus, ad calcem huius operis tabulam peculiarem adiungimus (tabulam II), quae pro valoribus ipsius $h$ a 0 usque ad 0,6 logarithmos respondentes ipsius $yy$ ad septem figuras decimales summa cura computatos exhibet. Argumentum $h$ a 0 usque ad 0,04 per singulas partes decies millesimas progreditur, quo pacto differentiae secundae ipsius $\log yy$ \SourceLeaf{127}evanescentes sunt redditae, ita ut in hac quidem tabulae parte interpolatio simplex sufficiat. Quoniam vero tabula, si ubivis eadem extensione gauderet, valde voluminosa evasisset, ab $h=0{,}04$ usque ad finem per singulas tantum millesimas partes progredi debuit; quamobrem in hac parte posteriori ad differentias secundas respicere oportebit, siquidem errores aliquot unitatum in figura septima evitare cupimus. Ceterum valores minores ipsius $h$ in praxi longe sunt frequentissimi.

Solutio aequationis 15 quoties $h$ limitem tabulae egreditur, nec non solutio aequationis $15^{\ast}$, sine difficultate per methodum indirectam vel per alias methodos satis cognitas perfici poterit. Ceterum haud abs re erit monere, valorem parvum ipsius $g$ cum valore negativo ipsius $\cos f$ consistere non posse nisi in orbitis valde excentricis, ut ex aequatione 20 infra in art. 95 tradenda sponte elucebit\footnote{Ostendit ista aequatio, si $\cos f$ sit negativus, $\varphi$ certe maiorem esse debere quam $90^0-g$.}.
